%%%%%%%%%%%%%%%%%%%%%%%%%%%%%%%%%%%%%%%%%%%%%%%%%%%%%%%%%%%%%%%%%%%%%%%%%%%%%%%%
% Rigorous Mathematics for Bilevel Geometry Learning with Conformal Calibration
% Option 2: Tuning-based approximation of the profiled (oracle) gradient
%%%%%%%%%%%%%%%%%%%%%%%%%%%%%%%%%%%%%%%%%%%%%%%%%%%%%%%%%%%%%%%%%%%%%%%%%%%%%%%%

\documentclass[11pt]{article}

\usepackage[margin=1in]{geometry}
\usepackage{amsmath,amssymb,amsthm,mathtools}
\usepackage{bm}
\usepackage{enumitem}
\usepackage{microtype}
\usepackage{algorithm}
\usepackage{algorithmic}

%%%%%%%%%%%%%%%%%%%%%%%%%%%%%%%%%%%%%%%%%%%%%%%%%%%%%%%%%%%%%%%%%%%%%%%%%%%%%%%%
% Notation
%%%%%%%%%%%%%%%%%%%%%%%%%%%%%%%%%%%%%%%%%%%%%%%%%%%%%%%%%%%%%%%%%%%%%%%%%%%%%%%%

\newcommand{\R}{\mathbb{R}}
\newcommand{\E}{\mathbb{E}}
\newcommand{\Prob}{\mathbb{P}}

\newcommand{\cX}{\mathcal{X}}
\newcommand{\cU}{\mathcal{U}}
\newcommand{\cB}{\mathcal{B}}
\newcommand{\cD}{\mathcal{D}}
\newcommand{\cS}{\mathcal{S}}

\newcommand{\ip}[2]{\left\langle #1, #2 \right\rangle}
\newcommand{\norm}[1]{\left\|#1\right\|}
\newcommand{\abs}[1]{\left|#1\right|}

\newcommand{\btheta}{\bm{\theta}}
\newcommand{\bu}{\bm{u}}
\newcommand{\bw}{\bm{w}}
\newcommand{\bx}{\bm{x}}
\newcommand{\bmu}{\bm{\mu}}

\DeclareMathOperator{\argmin}{argmin}
\DeclareMathOperator{\argmax}{argmax}
\DeclareMathOperator{\ri}{ri}

%%%%%%%%%%%%%%%%%%%%%%%%%%%%%%%%%%%%%%%%%%%%%%%%%%%%%%%%%%%%%%%%%%%%%%%%%%%%%%%%
% Theorem environments
%%%%%%%%%%%%%%%%%%%%%%%%%%%%%%%%%%%%%%%%%%%%%%%%%%%%%%%%%%%%%%%%%%%%%%%%%%%%%%%%

\theoremstyle{plain}
\newtheorem{theorem}{Theorem}
\newtheorem{lemma}[theorem]{Lemma}
\newtheorem{proposition}[theorem]{Proposition}
\newtheorem{corollary}[theorem]{Corollary}

\theoremstyle{definition}
\newtheorem{definition}[theorem]{Definition}
\newtheorem{assumption}[theorem]{Assumption}
\newtheorem{remark}[theorem]{Remark}

%%%%%%%%%%%%%%%%%%%%%%%%%%%%%%%%%%%%%%%%%%%%%%%%%%%%%%%%%%%%%%%%%%%%%%%%%%%%%%%%
% Document
%%%%%%%%%%%%%%%%%%%%%%%%%%%%%%%%%%%%%%%%%%%%%%%%%%%%%%%%%%%%%%%%%%%%%%%%%%%%%%%%

\title{Rigorous Mathematics for Robust Uncertainty-Set Geometry Learning\\
with Conformal Calibration (Option 2: Tuning-Based Profiled Gradients)}
\author{}
\date{}

\begin{document}
\maketitle

\begin{abstract}
We study robust optimization where the uncertainty distribution $P$ is unknown and the uncertainty set
$\cU_{\btheta,\rho}=\{\bu: s_{\btheta}(\bu)\le \rho\}$ is parameterized by a shape $\btheta$ and a size $\rho$.
The outer objective is to minimize the robust value $V(\btheta,\rho)$ subject to a coverage constraint
$\Prob(\bu\in \cU_{\btheta,\rho})\ge \tau$.
We present a rigorous, two-case mathematical development:
(i) when $P$ is known, the bilevel problem reduces to a profiled objective $J_P(\btheta)=V(\btheta,\rho_P(\btheta))$ and we derive the \emph{true} gradient of $J_P$,
including the quantile-sensitivity correction term that depends on $P$; (ii) when $P$ is unknown, we use split conformal calibration on an independent calibration set to obtain distribution-free marginal coverage.
We then formalize Option 2: use an additional \emph{tuning split} to construct a smooth, differentiable approximation of the profiled gradient (including the quantile-sensitivity term) and prove consistency of this approximate gradient under standard regularity conditions, while preserving final conformal validity via an independent calibration split.
\end{abstract}

\section{Setup}

\subsection{Parameterized uncertainty sets and gauge scores}

Fix an uncertainty dimension $d\ge 1$ and a parameter set $\Theta\subseteq \R^p$ (shape space).
We consider uncertainty sets of the form
\begin{equation}\label{eq:uset}
\cU_{\btheta,\rho} \;:=\; \{\bu\in \R^d:\; s_{\btheta}(\bu)\le \rho\},
\qquad \btheta\in\Theta,\;\rho>0,
\end{equation}
where $s_{\btheta}:\R^d\to [0,\infty)$ is the gauge (Minkowski functional) of the \emph{unit} set $\cU_{\btheta,1}$.

\begin{assumption}[Gauge regularity and scaling]\label{asm:gauge}
For each $\btheta\in\Theta$, the unit set $\cU_{\btheta,1}$ is nonempty, convex, compact, contains $\bm{0}$ in its interior, and
$s_{\btheta}$ is the gauge of $\cU_{\btheta,1}$:
\[
s_{\btheta}(\bu) \;=\; \inf\{t>0:\; \bu\in t\,\cU_{\btheta,1}\}.
\]
Consequently, $\cU_{\btheta,\rho}=\rho\,\cU_{\btheta,1}$ for all $\rho>0$.
\end{assumption}

\begin{definition}[Support function]
For a nonempty set $C\subseteq \R^d$, its support function is
\[
\sigma_C(\bw) \;:=\; \sup_{\bu\in C}\ip{\bw}{\bu} \in (-\infty,+\infty].
\]
\end{definition}

A standard convex-analytic identity is the scaling of support functions.

\begin{lemma}[Support function scaling]\label{lem:support-scaling}
If $C\subseteq \R^d$ is nonempty and $\rho>0$, then
\[
\sigma_{\rho C}(\bw) = \rho\,\sigma_C(\bw)\qquad \forall \bw\in \R^d.
\]
\end{lemma}

\begin{proof}
By definition,
\[
\sigma_{\rho C}(\bw)
=\sup_{\bu\in \rho C}\ip{\bw}{\bu}
=\sup_{\bv\in C}\ip{\bw}{\rho \bv}
=\rho \sup_{\bv\in C}\ip{\bw}{\bv}
=\rho\,\sigma_C(\bw).
\]
\end{proof}

Under Assumption~\ref{asm:gauge}, Lemma~\ref{lem:support-scaling} yields
\begin{equation}\label{eq:sigma-scale}
\sigma_{\cU_{\btheta,\rho}}(\bw) = \rho\,\sigma_{\cU_{\btheta,1}}(\bw).
\end{equation}

\subsection{Robust optimization layer and value function}

Let $\cX\subseteq \R^{n_x}$ be a closed convex decision set, $f:\R^{n_x}\to \R\cup\{+\infty\}$ be a proper closed convex objective,
and $a_j:\R^{n_x}\to \R\cup\{+\infty\}$ be proper closed convex functions.
Fix exposure vectors $\bw_1,\dots,\bw_m\in \R^d$ and scalars $b_1,\dots,b_m\in \R$.

Given $(\btheta,\rho)$, define the robust optimization problem
\begin{equation}\label{eq:robust}
\begin{aligned}
V(\btheta,\rho)
\;:=\;
\inf_{\bx\in \cX}\quad & f(\bx)\\
\text{s.t.}\quad &
a_j(\bx)+\sigma_{\cU_{\btheta,\rho}}(\bw_j)\le b_j,\qquad j=1,\dots,m.
\end{aligned}
\end{equation}

\begin{assumption}[Convexity and Slater]\label{asm:slater}
For each $(\btheta,\rho)$:
(i) $\cX$ is nonempty, closed, convex; (ii) $f$ and all $a_j$ are proper closed convex; and
(iii) there exists $\bar{\bx}\in \ri(\cX)$ such that
$a_j(\bar{\bx})+\sigma_{\cU_{\btheta,\rho}}(\bw_j) < b_j$ for all $j$ (Slater condition).
\end{assumption}

Assumption~\ref{asm:slater} ensures strong duality and existence of optimal dual multipliers for~\eqref{eq:robust}.

\subsection{Outer problem and score quantiles}

Let $U\sim P$ be a random uncertainty realization in $\R^d$.
Define the gauge score random variable
\[
S_{\btheta} := s_{\btheta}(U).
\]
Define its CDF (under $P$)
\[
F_{\btheta}(r) := \Prob(S_{\btheta}\le r),\qquad r\in \R.
\]
Fix a target coverage level $\tau\in(0,1)$.

The learning problem is:
\begin{equation}\label{eq:bilevel}
\boxed{
\min_{\btheta\in \Theta,\;\rho>0} V(\btheta,\rho)
\quad\text{s.t.}\quad
\Prob(U\in \cU_{\btheta,\rho})\ge \tau.
}
\end{equation}
Since $U\in \cU_{\btheta,\rho}$ iff $s_{\btheta}(U)\le \rho$, the coverage constraint is
\begin{equation}\label{eq:coverage}
\Prob(U\in \cU_{\btheta,\rho})\ge \tau
\quad\Longleftrightarrow\quad
F_{\btheta}(\rho)\ge \tau.
\end{equation}

We define the (left) $\tau$-quantile radius map:
\begin{equation}\label{eq:rhoP}
\rho_P(\btheta)\;:=\;\inf\{r>0:\;F_{\btheta}(r)\ge \tau\}\in (0,+\infty].
\end{equation}

\section{Case I: $P$ known (oracle profiling and the true gradient)}

\subsection{Profiling out the radius}

\begin{assumption}[Monotonicity in radius]\label{asm:monotone}
For each $\btheta\in\Theta$ and $0<\rho_1\le \rho_2$, $\cU_{\btheta,\rho_1}\subseteq \cU_{\btheta,\rho_2}$.
\end{assumption}

Assumption~\ref{asm:monotone} holds under Assumption~\ref{asm:gauge} by construction.

\begin{proposition}[Radius profiling]\label{prop:profiling}
Assume Assumptions~\ref{asm:gauge} and~\ref{asm:slater}. Fix $\btheta\in\Theta$ and suppose $\rho_P(\btheta)<\infty$.
Then:
\begin{enumerate}[label=(\roman*), leftmargin=2em, itemsep=1pt]
\item The map $\rho\mapsto V(\btheta,\rho)$ is nondecreasing on $(0,\infty)$.
\item The smallest feasible radius for the coverage constraint is $\rho=\rho_P(\btheta)$.
\item The bilevel problem~\eqref{eq:bilevel} is equivalent to the single-level \emph{profiled} problem
\begin{equation}\label{eq:profiled}
\min_{\btheta\in\Theta}\; J_P(\btheta)
\qquad\text{where}\qquad
J_P(\btheta):=V\big(\btheta,\rho_P(\btheta)\big).
\end{equation}
\end{enumerate}
\end{proposition}

\begin{proof}
(i) Fix $\btheta$ and let $0<\rho_1\le \rho_2$. By Assumption~\ref{asm:monotone},
$\cU_{\btheta,\rho_1}\subseteq \cU_{\btheta,\rho_2}$, hence for each $j$,
\[
\sigma_{\cU_{\btheta,\rho_1}}(\bw_j)
=\sup_{\bu\in \cU_{\btheta,\rho_1}}\ip{\bw_j}{\bu}
\le
\sup_{\bu\in \cU_{\btheta,\rho_2}}\ip{\bw_j}{\bu}
=
\sigma_{\cU_{\btheta,\rho_2}}(\bw_j).
\]
Therefore, every feasible $\bx$ for radius $\rho_2$ is also feasible for radius $\rho_1$ (the constraints at $\rho_1$ are weaker),
so the feasible region of~\eqref{eq:robust} expands as $\rho$ decreases. Thus the optimal value is nondecreasing in $\rho$:
$V(\btheta,\rho_1)\le V(\btheta,\rho_2)$.

(ii) By~\eqref{eq:coverage} and~\eqref{eq:rhoP}, the coverage constraint $F_{\btheta}(\rho)\ge \tau$ holds iff $\rho\ge \rho_P(\btheta)$.
Hence the smallest feasible radius is $\rho_P(\btheta)$.

(iii) Fix any feasible $(\btheta,\rho)$ for~\eqref{eq:bilevel}. By (ii), $\rho\ge \rho_P(\btheta)$.
By (i), $V(\btheta,\rho)\ge V(\btheta,\rho_P(\btheta))=J_P(\btheta)$.
Thus, for each $\btheta$, the best feasible choice is $\rho=\rho_P(\btheta)$, and the bilevel reduces to~\eqref{eq:profiled}.
\end{proof}

\subsection{Envelope derivatives of $V(\btheta,\rho)$}

We now derive derivatives of the robust value function with respect to $(\btheta,\rho)$ without differentiating through the optimizer.

Define $\sigma_j(\btheta,\rho):=\sigma_{\cU_{\btheta,\rho}}(\bw_j)$.

\begin{assumption}[Differentiability of support functions]\label{asm:sigma-diff}
Fix $\rho>0$. For each $j\in\{1,\dots,m\}$, the map $\btheta\mapsto \sigma_j(\btheta,\rho)$ is differentiable on $\Theta$.
\end{assumption}

\begin{theorem}[Envelope gradient for shape]\label{thm:envelope-theta}
Assume Assumptions~\ref{asm:slater} and~\ref{asm:sigma-diff}. Fix $(\btheta,\rho)$.
Let $\bmu^*(\btheta,\rho)\in \R_+^m$ be any optimal dual multiplier vector for~\eqref{eq:robust}.
Then
\begin{equation}\label{eq:env-theta}
\bg_{\btheta}(\btheta,\rho)
:=\sum_{j=1}^m \mu_j^*(\btheta,\rho)\,\nabla_{\btheta}\sigma_{\cU_{\btheta,\rho}}(\bw_j)
\end{equation}
is a subgradient of $\btheta\mapsto V(\btheta,\rho)$ at $\btheta$.
If, additionally, the optimal dual multiplier is unique at $(\btheta,\rho)$, then $V(\cdot,\rho)$ is differentiable at $\btheta$ and
\[
\nabla_{\btheta}V(\btheta,\rho)=\bg_{\btheta}(\btheta,\rho).
\]
\end{theorem}

\begin{proof}
We write the Lagrangian of~\eqref{eq:robust}. Introduce the indicator $I_{\cX}(\bx)$ which is $0$ if $\bx\in\cX$ and $+\infty$ otherwise.
Define constraint functions
\[
g_j(\bx;\btheta,\rho) := a_j(\bx) + \sigma_j(\btheta,\rho) - b_j.
\]
Then~\eqref{eq:robust} is $\inf_{\bx} f(\bx)+I_{\cX}(\bx)$ subject to $g_j(\bx;\btheta,\rho)\le 0$.

For multipliers $\bmu\in \R_+^m$, the Lagrangian is
\[
\cL(\bx,\bmu;\btheta,\rho)
=
f(\bx)+I_{\cX}(\bx)+\sum_{j=1}^m \mu_j\big(a_j(\bx)+\sigma_j(\btheta,\rho)-b_j\big).
\]
Define the dual function
\[
q(\bmu;\btheta,\rho) := \inf_{\bx\in \R^{n_x}} \cL(\bx,\bmu;\btheta,\rho).
\]
Because $\sigma_j(\btheta,\rho)$ does not depend on $\bx$, we can separate:
\[
q(\bmu;\btheta,\rho)
=
\underbrace{\inf_{\bx}\left(f(\bx)+I_{\cX}(\bx)+\sum_{j=1}^m \mu_j a_j(\bx)\right)}_{=: \tilde q(\bmu)\ \text{(independent of $\btheta,\rho$)}}
\;+\;
\sum_{j=1}^m \mu_j\big(\sigma_j(\btheta,\rho)-b_j\big).
\]
Hence, for fixed $(\rho,\bmu)$, $q(\bmu;\cdot,\rho)$ is differentiable with
\[
\nabla_{\btheta} q(\bmu;\btheta,\rho) = \sum_{j=1}^m \mu_j\,\nabla_{\btheta}\sigma_j(\btheta,\rho).
\]

By Assumption~\ref{asm:slater}, strong duality holds and the dual optimum is attained:
\[
V(\btheta,\rho) = \max_{\bmu\ge 0} q(\bmu;\btheta,\rho),
\qquad
\bmu^*(\btheta,\rho)\in \argmax_{\bmu\ge 0} q(\bmu;\btheta,\rho).
\]
We now prove that $\bg_{\btheta}(\btheta,\rho)$ in~\eqref{eq:env-theta} is a subgradient.
Fix $\btheta'$ and let $\bmu^*=\bmu^*(\btheta,\rho)$ be an optimal multiplier at $(\btheta,\rho)$.
Then
\[
V(\btheta',\rho) = \max_{\bmu\ge 0} q(\bmu;\btheta',\rho) \;\ge\; q(\bmu^*;\btheta',\rho).
\]
By differentiability of $q(\bmu^*;\cdot,\rho)$ at $\btheta$,
\[
q(\bmu^*;\btheta',\rho)
=
q(\bmu^*;\btheta,\rho)
+
\ip{\nabla_{\btheta}q(\bmu^*;\btheta,\rho)}{\btheta'-\btheta}
+o(\norm{\btheta'-\btheta}).
\]
Using $q(\bmu^*;\btheta,\rho)=V(\btheta,\rho)$ (optimality),
\[
V(\btheta',\rho)\ge V(\btheta,\rho) + \ip{\nabla_{\btheta}q(\bmu^*;\btheta,\rho)}{\btheta'-\btheta} + o(\norm{\btheta'-\btheta}).
\]
Thus $\nabla_{\btheta}q(\bmu^*;\btheta,\rho)$ is a subgradient of $V(\cdot,\rho)$ at $\btheta$.
Finally,
\[
\nabla_{\btheta}q(\bmu^*;\btheta,\rho)=\sum_{j=1}^m \mu_j^*(\btheta,\rho)\,\nabla_{\btheta}\sigma_j(\btheta,\rho)=\bg_{\btheta}(\btheta,\rho).
\]
If the dual maximizer is unique, standard Danskin/envelope arguments imply differentiability of $V(\cdot,\rho)$ at $\btheta$ and the subgradient is the gradient.
\end{proof}

Next, we obtain an envelope derivative with respect to $\rho$.

\begin{theorem}[Envelope derivative for size]\label{thm:envelope-rho}
Assume Assumptions~\ref{asm:gauge} and~\ref{asm:slater}. Fix $(\btheta,\rho)$.
Let $\bmu^*(\btheta,\rho)$ be any optimal dual multipliers for~\eqref{eq:robust}.
Then
\begin{equation}\label{eq:env-rho}
g_{\rho}(\btheta,\rho)
:=
\sum_{j=1}^m \mu_j^*(\btheta,\rho)\,\sigma_{\cU_{\btheta,1}}(\bw_j)
\end{equation}
is a subgradient of $\rho\mapsto V(\btheta,\rho)$ at $\rho$.
If the optimal dual multiplier is unique at $(\btheta,\rho)$, then $V(\btheta,\cdot)$ is differentiable at $\rho$ and
\[
\frac{\partial}{\partial\rho}V(\btheta,\rho)=g_\rho(\btheta,\rho).
\]
\end{theorem}

\begin{proof}
From the dual representation in the proof of Theorem~\ref{thm:envelope-theta},
\[
V(\btheta,\rho)=\max_{\bmu\ge 0}\Big(\tilde q(\bmu)+\sum_{j=1}^m \mu_j(\sigma_j(\btheta,\rho)-b_j)\Big).
\]
By Lemma~\ref{lem:support-scaling} and Assumption~\ref{asm:gauge},
$\sigma_j(\btheta,\rho)=\rho\,\sigma_{\cU_{\btheta,1}}(\bw_j)$, hence for fixed $(\btheta,\bmu)$ the dual function is affine in $\rho$:
\[
q(\bmu;\btheta,\rho)=
\tilde q(\bmu)-\sum_{j=1}^m \mu_j b_j
+\rho \sum_{j=1}^m \mu_j \sigma_{\cU_{\btheta,1}}(\bw_j).
\]
Therefore, for each fixed $(\btheta,\bmu)$, $\frac{\partial}{\partial\rho}q(\bmu;\btheta,\rho)=\sum_j \mu_j \sigma_{\cU_{\btheta,1}}(\bw_j)$.
Now repeat the same subgradient argument as in Theorem~\ref{thm:envelope-theta}, treating $\rho$ as the parameter and using any optimal maximizer $\bmu^*(\btheta,\rho)$.
This gives~\eqref{eq:env-rho}. Uniqueness of $\bmu^*$ yields differentiability and equality.
\end{proof}

\subsection{The true gradient of the oracle profiled objective}

By Proposition~\ref{prop:profiling}, the oracle profiled objective is
\[
J_P(\btheta)=V\big(\btheta,\rho_P(\btheta)\big).
\]
The \emph{true} gradient of $J_P$ includes both the direct effect of $\btheta$ on $V$ and the indirect effect through $\rho_P(\btheta)$.

To state a differentiability result for $\rho_P(\btheta)$, we impose a standard regularity condition at the quantile.

\begin{assumption}[Quantile regularity]\label{asm:quantile-reg}
Fix $\btheta\in\Theta$ and let $\rho_P(\btheta)$ be defined as in~\eqref{eq:rhoP}.
There exists an open interval $I$ containing $\rho_P(\btheta)$ such that:
(i) $F_{\btheta}$ is continuously differentiable on $I$ with density $f_{\btheta}(r):=\partial_r F_{\btheta}(r)$;
(ii) $f_{\btheta}(\rho_P(\btheta))>0$; and
(iii) for $\btheta'$ in a neighborhood of $\btheta$, the map $(\btheta',r)\mapsto F_{\btheta'}(r)$ is continuously differentiable on $\{\btheta'\}\times I$.
\end{assumption}

\begin{lemma}[Implicit derivative of the quantile map]\label{lem:quantile-derivative}
Assume Assumption~\ref{asm:quantile-reg} holds at $\btheta$.
Then there exists a neighborhood $N$ of $\btheta$ and a continuously differentiable map $\rho_P:N\to \R$ such that
$F_{\btheta'}(\rho_P(\btheta'))=\tau$ for all $\btheta'\in N$.
Moreover,
\begin{equation}\label{eq:quantile-ift}
\nabla_{\btheta}\rho_P(\btheta)
=
-\frac{\nabla_{\btheta}F_{\btheta}(\rho_P(\btheta))}{f_{\btheta}(\rho_P(\btheta))}.
\end{equation}
\end{lemma}

\begin{proof}
Define $G(\btheta,r):=F_{\btheta}(r)-\tau$.
By definition of $\rho_P(\btheta)$ under Assumption~\ref{asm:quantile-reg}, $G(\btheta,\rho_P(\btheta))=0$.
Also, $\partial_r G(\btheta,\rho_P(\btheta))=f_{\btheta}(\rho_P(\btheta))>0$ by Assumption~\ref{asm:quantile-reg}(ii).
By the implicit function theorem, there exists a neighborhood $N$ of $\btheta$ and a unique continuously differentiable function $\rho_P$ on $N$
such that $G(\btheta',\rho_P(\btheta'))=0$ for all $\btheta'\in N$, i.e., $F_{\btheta'}(\rho_P(\btheta'))=\tau$.
Differentiating $G(\btheta,\rho_P(\btheta))=0$ with respect to $\btheta$ yields
\[
\nabla_{\btheta}G(\btheta,\rho_P(\btheta)) + \partial_r G(\btheta,\rho_P(\btheta)) \,\nabla_{\btheta}\rho_P(\btheta)=0,
\]
which is exactly~\eqref{eq:quantile-ift}.
\end{proof}

\begin{theorem}[True oracle gradient of the profiled objective]\label{thm:true-oracle-grad}
Assume:
(i) Assumptions~\ref{asm:gauge},~\ref{asm:slater}, and~\ref{asm:sigma-diff} hold;
(ii) Assumption~\ref{asm:quantile-reg} holds at $\btheta$; and
(iii) $V(\cdot,\rho)$ is differentiable at $\btheta$ for $\rho=\rho_P(\btheta)$ and $V(\btheta,\cdot)$ is differentiable at $\rho_P(\btheta)$.
Then $J_P$ is differentiable at $\btheta$ and
\begin{equation}\label{eq:true-grad}
\boxed{
\nabla_{\btheta}J_P(\btheta)
=
\nabla_{\btheta}V\big(\btheta,\rho_P(\btheta)\big)
+
\frac{\partial}{\partial\rho}V\big(\btheta,\rho_P(\btheta)\big)\,\nabla_{\btheta}\rho_P(\btheta).
}
\end{equation}
Moreover, letting $\bmu^*=\bmu^*(\btheta,\rho_P(\btheta))$ be the (unique) optimal dual multipliers for~\eqref{eq:robust} at $(\btheta,\rho_P(\btheta))$,
\begin{equation}\label{eq:true-grad-expanded}
\boxed{
\nabla_{\btheta}J_P(\btheta)
=
\sum_{j=1}^m \mu_j^*\,\nabla_{\btheta}\sigma_{\cU_{\btheta,\rho_P(\btheta)}}(\bw_j)
+
\left(\sum_{j=1}^m \mu_j^*\,\sigma_{\cU_{\btheta,1}}(\bw_j)\right)
\left(
-\frac{\nabla_{\btheta}F_{\btheta}(\rho_P(\btheta))}{f_{\btheta}(\rho_P(\btheta))}
\right).
}
\end{equation}
\end{theorem}

\begin{proof}
The chain rule gives~\eqref{eq:true-grad} because $J_P(\btheta)=V(\btheta,\rho_P(\btheta))$ and both factors are differentiable by assumption.
The first term in~\eqref{eq:true-grad-expanded} follows from Theorem~\ref{thm:envelope-theta} (uniqueness gives equality).
The scalar derivative $\partial_\rho V$ is given by Theorem~\ref{thm:envelope-rho}.
Finally, $\nabla_{\btheta}\rho_P(\btheta)$ is given by Lemma~\ref{lem:quantile-derivative}.
Substituting these identities into~\eqref{eq:true-grad} yields~\eqref{eq:true-grad-expanded}.
\end{proof}

\begin{remark}[Where $P$ enters the gradient]\label{rem:P-enters}
The second term in~\eqref{eq:true-grad} is the \emph{quantile-sensitivity correction}.
It depends on $P$ through $\nabla_{\btheta}F_{\btheta}(\rho_P(\btheta))$ and $f_{\btheta}(\rho_P(\btheta))$.
If $P$ is unknown and no parametric model is assumed, this correction cannot be evaluated exactly; this is the fundamental reason that naive ``optimize $\btheta$ then calibrate $\rho$'' is not literally optimizing the oracle profiled objective $J_P$.
Option 2 addresses this by approximating the correction on a tuning split.
\end{remark}

\section{Case II: $P$ unknown (distribution-free conformal calibration)}

We now assume $P$ is unknown and only sample access is available.
Let $\{U_i\}_{i=1}^{n_{\mathrm{cal}}}$ be a calibration sample (typically held out from training/tuning).
We assume minimal symmetry:

\begin{assumption}[Exchangeability for conformal]\label{asm:exchange}
$(U_1,\dots,U_{n_{\mathrm{cal}}},U_{\mathrm{new}})$ are exchangeable.
\end{assumption}

\begin{theorem}[Split conformal marginal coverage]\label{thm:conformal}
Assume Assumption~\ref{asm:exchange}. Fix any (possibly random) shape $\btheta$ that is independent of the calibration sample $(U_i)_{i=1}^{n_{\mathrm{cal}}}$.
Define calibration scores $S_i=s_{\btheta}(U_i)$ for $i=1,\dots,n_{\mathrm{cal}}$ and let
$S_{(1)}\le \dots\le S_{(n_{\mathrm{cal}})}$ be their order statistics.
Set $S_{(n_{\mathrm{cal}}+1)}:=+\infty$ and define
\begin{equation}\label{eq:conformal-rho}
k:=\lceil (n_{\mathrm{cal}}+1)\tau\rceil,\qquad
\hat\rho_{\tau}:=S_{(k)}.
\end{equation}
Then
\[
\Prob\big(s_{\btheta}(U_{\mathrm{new}})\le \hat\rho_\tau\big)\ge \tau,
\qquad\text{equivalently}\qquad
\Prob\big(U_{\mathrm{new}}\in \cU_{\btheta,\hat\rho_\tau}\big)\ge \tau.
\]
\end{theorem}

\begin{proof}
Consider the $(n_{\mathrm{cal}}+1)$ scores
\[
S_1,\dots,S_{n_{\mathrm{cal}}},S_{\mathrm{new}}
\quad\text{where}\quad
S_{\mathrm{new}}:=s_{\btheta}(U_{\mathrm{new}}).
\]
By Assumption~\ref{asm:exchange} and independence of $\btheta$ from the calibration sample,
the joint distribution of these scores is exchangeable.
Therefore, the rank of $S_{\mathrm{new}}$ among $\{S_1,\dots,S_{n_{\mathrm{cal}}},S_{\mathrm{new}}\}$ is uniform on $\{1,\dots,n_{\mathrm{cal}}+1\}$ (ties can be handled by the $+\infty$ augmentation).
The event $S_{\mathrm{new}}\le S_{(k)}$ occurs whenever the rank of $S_{\mathrm{new}}$ is at most $k$, hence
\[
\Prob(S_{\mathrm{new}}\le S_{(k)})
=
\Prob(\mathrm{rank}(S_{\mathrm{new}})\le k)
=
\frac{k}{n_{\mathrm{cal}}+1}
\ge \tau,
\]
since $k=\lceil (n_{\mathrm{cal}}+1)\tau\rceil$.
\end{proof}

\begin{remark}[Independence requirement]\label{rem:indep}
Theorem~\ref{thm:conformal} only requires that the final calibrated radius $\hat\rho_\tau$ be computed on data independent of how $\btheta$ was obtained.
In Option 2 below, $\btheta$ may depend on a separate training split and a separate tuning split; final calibration is performed on an independent calibration split, so the conformal guarantee remains valid.
\end{remark}

\section{Option 2: Tuning-based approximation of the profiled gradient}

This section formalizes Option 2: approximate the oracle profiled gradient in Case I using an auxiliary tuning split, while preserving the final distribution-free coverage guarantee via an independent calibration split.

\subsection{Data splits}

Assume we have i.i.d.\ data $U_1,\dots,U_N\sim P$ (or exchangeable, where needed).
We partition indices into three disjoint sets:
\[
\cD_{\mathrm{tr}}=\{U_i: i\in I_{\mathrm{tr}}\},\qquad
\cD_{\mathrm{tune}}=\{U_i: i\in I_{\mathrm{tune}}\},\qquad
\cD_{\mathrm{cal}}=\{U_i: i\in I_{\mathrm{cal}}\},
\]
with sizes $n_{\mathrm{tr}}, n_{\mathrm{tune}}, n_{\mathrm{cal}}$.

\begin{itemize}[leftmargin=2em, itemsep=1pt]
\item The \textbf{training} split $\cD_{\mathrm{tr}}$ is used to optimize (or approximate) gradients of the robust objective.
\item The \textbf{tuning} split $\cD_{\mathrm{tune}}$ is used to approximate the \emph{quantile sensitivity} $\nabla_{\btheta}\rho_P(\btheta)$.
\item The \textbf{calibration} split $\cD_{\mathrm{cal}}$ is used \emph{only} for conformal radius calibration, guaranteeing marginal coverage.
\end{itemize}

\subsection{A smooth quantile map that is differentiable in $\btheta$}

The main technical obstacle in~\eqref{eq:true-grad-expanded} is that $\btheta\mapsto \rho_P(\btheta)$ is defined through a nonsmooth indicator CDF $F_{\btheta}(r)=\E[\mathbf{1}\{s_{\btheta}(U)\le r\}]$.
We introduce a smooth approximation (standard in implicit differentiation for quantiles).

\begin{assumption}[Smoothing function]\label{asm:smoothing}
Let $\Phi:\R\to[0,1]$ be continuously differentiable, nondecreasing, and satisfy
$\lim_{t\to-\infty}\Phi(t)=0$ and $\lim_{t\to+\infty}\Phi(t)=1$.
Let $\varphi(t):=\Phi'(t)$ and assume $\varphi$ is bounded and not identically zero.
\end{assumption}

Fix a smoothing bandwidth $\varepsilon>0$.
Define the smoothed CDF
\begin{equation}\label{eq:smoothed-cdf}
F_{\btheta,\varepsilon}(r)
:=
\E\left[\Phi\left(\frac{r-s_{\btheta}(U)}{\varepsilon}\right)\right].
\end{equation}
This is a continuous, differentiable function of $(\btheta,r)$ whenever $\btheta\mapsto s_{\btheta}(u)$ is differentiable.

\begin{assumption}[Score differentiability and domination]\label{asm:score-diff}
For each $u\in\R^d$, the map $\btheta\mapsto s_{\btheta}(u)$ is continuously differentiable.
Moreover, there exists an integrable function $M(U)$ such that $\norm{\nabla_{\btheta}s_{\btheta}(U)}\le M(U)$ almost surely for all $\btheta$ in a neighborhood of interest.
\end{assumption}

\begin{lemma}[Derivatives of the smoothed CDF]\label{lem:smoothed-derivatives}
Assume Assumptions~\ref{asm:smoothing} and~\ref{asm:score-diff}. Then for each $\varepsilon>0$:
\begin{align}
\partial_r F_{\btheta,\varepsilon}(r)
&=
\frac{1}{\varepsilon}\E\left[\varphi\left(\frac{r-s_{\btheta}(U)}{\varepsilon}\right)\right],\label{eq:dFr}\\
\nabla_{\btheta}F_{\btheta,\varepsilon}(r)
&=
-\frac{1}{\varepsilon}\E\left[\varphi\left(\frac{r-s_{\btheta}(U)}{\varepsilon}\right)\,\nabla_{\btheta}s_{\btheta}(U)\right].\label{eq:dFtheta}
\end{align}
\end{lemma}

\begin{proof}
The integrand in~\eqref{eq:smoothed-cdf} is differentiable in $(r,\btheta)$ for each realization $U$.
Assumption~\ref{asm:smoothing} gives bounded $\varphi$, and Assumption~\ref{asm:score-diff} provides an integrable envelope for $\nabla_{\btheta}s_{\btheta}(U)$.
Thus differentiation under the expectation is justified by dominated convergence.
Applying the chain rule yields~\eqref{eq:dFr} and~\eqref{eq:dFtheta}.
\end{proof}

Define the \emph{smoothed $\tau$-quantile radius}:
\begin{equation}\label{eq:rho-eps}
\rho_{P,\varepsilon}(\btheta)
\;:=\;
\inf\{r:\;F_{\btheta,\varepsilon}(r)\ge \tau\}.
\end{equation}

To ensure uniqueness and differentiability, we assume strict monotonicity at the solution.

\begin{assumption}[Smoothed quantile regularity]\label{asm:smoothed-quantile-reg}
Fix $\btheta$ and $\varepsilon>0$. Assume there exists $\rho_{P,\varepsilon}(\btheta)\in\R$ such that
$F_{\btheta,\varepsilon}(\rho_{P,\varepsilon}(\btheta))=\tau$ and
$\partial_r F_{\btheta,\varepsilon}(\rho_{P,\varepsilon}(\btheta))>0$.
\end{assumption}

\begin{lemma}[Implicit derivative of the smoothed quantile map]\label{lem:smoothed-quantile-deriv}
Assume Assumptions~\ref{asm:smoothing},~\ref{asm:score-diff}, and~\ref{asm:smoothed-quantile-reg}.
Then $\rho_{P,\varepsilon}$ is differentiable at $\btheta$ and
\begin{equation}\label{eq:rhoeps-deriv}
\nabla_{\btheta}\rho_{P,\varepsilon}(\btheta)
=
-\frac{\nabla_{\btheta}F_{\btheta,\varepsilon}(\rho_{P,\varepsilon}(\btheta))}{\partial_r F_{\btheta,\varepsilon}(\rho_{P,\varepsilon}(\btheta))}
=
\frac{
\E\left[\varphi\left(\frac{\rho_{P,\varepsilon}(\btheta)-s_{\btheta}(U)}{\varepsilon}\right)\nabla_{\btheta}s_{\btheta}(U)\right]
}{
\E\left[\varphi\left(\frac{\rho_{P,\varepsilon}(\btheta)-s_{\btheta}(U)}{\varepsilon}\right)\right]
}.
\end{equation}
\end{lemma}

\begin{proof}
The first equality is the implicit function theorem applied to $G(\btheta,r)=F_{\btheta,\varepsilon}(r)-\tau$,
using $\partial_r G=\partial_r F_{\btheta,\varepsilon}>0$ at the root.
The second equality follows by substituting Lemma~\ref{lem:smoothed-derivatives} into the ratio, noting that the factor $1/\varepsilon$ cancels.
\end{proof}

\subsection{Empirical tuning estimators}

Given a tuning sample $\cD_{\mathrm{tune}}=\{U_i\}_{i=1}^{n_{\mathrm{tune}}}$, define scores $S_i(\btheta)=s_{\btheta}(U_i)$ and the empirical smoothed CDF
\begin{equation}\label{eq:Fhat}
\hat F_{\btheta,\varepsilon}(r)
:=
\frac{1}{n_{\mathrm{tune}}}
\sum_{i=1}^{n_{\mathrm{tune}}}
\Phi\left(\frac{r-S_i(\btheta)}{\varepsilon}\right).
\end{equation}
Define the empirical smoothed quantile radius
\begin{equation}\label{eq:rhohat-eps}
\hat\rho_{\varepsilon}(\btheta)
\;:=\;
\inf\{r:\;\hat F_{\btheta,\varepsilon}(r)\ge \tau\}.
\end{equation}

For the gradient of $\hat\rho_{\varepsilon}(\btheta)$, we use the finite-sample analogue of~\eqref{eq:rhoeps-deriv}.
Let weights
\[
\omega_i(\btheta)
:=
\varphi\left(\frac{\hat\rho_{\varepsilon}(\btheta)-S_i(\btheta)}{\varepsilon}\right)\ge 0.
\]
Assuming $\sum_i \omega_i(\btheta)>0$, define
\begin{equation}\label{eq:rhohat-grad}
\widehat{\nabla_{\btheta}\rho_{\varepsilon}}(\btheta)
:=
\frac{\sum_{i=1}^{n_{\mathrm{tune}}}\omega_i(\btheta)\,\nabla_{\btheta}s_{\btheta}(U_i)}{\sum_{i=1}^{n_{\mathrm{tune}}}\omega_i(\btheta)}.
\end{equation}

\begin{remark}[Interpretation of~\eqref{eq:rhohat-grad}]
The weights $\omega_i(\btheta)$ concentrate on points with $S_i(\btheta)$ close to $\hat\rho_{\varepsilon}(\btheta)$.
Thus~\eqref{eq:rhohat-grad} is a soft analogue of a boundary conditional expectation:
$\nabla_{\btheta}\rho_P(\btheta)\approx \E[\nabla_{\btheta}s_{\btheta}(U)\mid s_{\btheta}(U)=\rho_P(\btheta)]$.
\end{remark}

\subsection{Approximate profiled gradient estimator}

Define the smoothed profiled objective:
\begin{equation}\label{eq:JPeps}
J_{P,\varepsilon}(\btheta)
:=
V\big(\btheta,\rho_{P,\varepsilon}(\btheta)\big).
\end{equation}
By the same chain rule as Theorem~\ref{thm:true-oracle-grad}, if differentiable,
\[
\nabla_{\btheta}J_{P,\varepsilon}(\btheta)
=
\nabla_{\btheta}V(\btheta,\rho_{P,\varepsilon}(\btheta))
+
\partial_{\rho}V(\btheta,\rho_{P,\varepsilon}(\btheta))\,
\nabla_{\btheta}\rho_{P,\varepsilon}(\btheta).
\]

We approximate $\rho_{P,\varepsilon}(\btheta)$ and $\nabla_{\btheta}\rho_{P,\varepsilon}(\btheta)$ via the tuning estimators
$\hat\rho_{\varepsilon}(\btheta)$ and $\widehat{\nabla_{\btheta}\rho_{\varepsilon}}(\btheta)$,
and approximate the envelope terms $\nabla_{\btheta}V$ and $\partial_\rho V$ by solving the robust problem at $(\btheta,\hat\rho_{\varepsilon}(\btheta))$.

Let $\bmu^*(\btheta,\hat\rho_{\varepsilon}(\btheta))$ denote an optimal dual multiplier for the robust solve at $(\btheta,\hat\rho_{\varepsilon}(\btheta))$.
Define the approximate gradient estimator
\begin{equation}\label{eq:ghat}
\boxed{
\hat{\bg}_{\varepsilon}(\btheta)
:=
\underbrace{\sum_{j=1}^m \mu_j^*(\btheta,\hat\rho_{\varepsilon}(\btheta))\,
\nabla_{\btheta}\sigma_{\cU_{\btheta,\hat\rho_{\varepsilon}(\btheta)}}(\bw_j)}_{\text{envelope shape term}}
\;+\;
\underbrace{\left(\sum_{j=1}^m \mu_j^*(\btheta,\hat\rho_{\varepsilon}(\btheta))\,\sigma_{\cU_{\btheta,1}}(\bw_j)\right)}_{\text{envelope size sensitivity}}
\cdot
\underbrace{\widehat{\nabla_{\btheta}\rho_{\varepsilon}}(\btheta)}_{\text{tuned quantile sensitivity}}.
}
\end{equation}

\subsection{Consistency of the tuned gradient (fixed $\varepsilon$)}

We state a consistency theorem under standard regularity.
To avoid technicalities about set-valued dual solutions, we assume uniqueness/continuity of the dual multiplier map.

\begin{assumption}[Dual uniqueness and continuity]\label{asm:dual-unique}
For each $(\btheta,\rho)$ in a neighborhood of interest, the robust problem~\eqref{eq:robust} has a unique optimal dual multiplier vector $\bmu^*(\btheta,\rho)$.
Moreover, $(\btheta,\rho)\mapsto \bmu^*(\btheta,\rho)$ is continuous.
\end{assumption}

\begin{assumption}[Continuity of support sensitivities]\label{asm:sensitivity-cont}
For each $j$, the maps $(\btheta,\rho)\mapsto \nabla_{\btheta}\sigma_{\cU_{\btheta,\rho}}(\bw_j)$ and $\btheta\mapsto \sigma_{\cU_{\btheta,1}}(\bw_j)$
are continuous on the region of interest.
\end{assumption}

\begin{theorem}[Consistency of the tuned profiled gradient for fixed $\varepsilon$]\label{thm:consistency-eps}
Assume:
\begin{itemize}[leftmargin=2em, itemsep=1pt]
\item $U_1,\dots,U_{n_{\mathrm{tune}}}$ are i.i.d.\ from $P$;
\item Assumptions~\ref{asm:smoothing},~\ref{asm:score-diff}, and~\ref{asm:smoothed-quantile-reg} hold;
\item Assumptions~\ref{asm:slater},~\ref{asm:gauge},~\ref{asm:sigma-diff},~\ref{asm:dual-unique}, and~\ref{asm:sensitivity-cont} hold.
\end{itemize}
Fix $\varepsilon>0$ and a shape $\btheta$.
Then, as $n_{\mathrm{tune}}\to\infty$,
\[
\hat\rho_{\varepsilon}(\btheta)\to \rho_{P,\varepsilon}(\btheta)
\quad\text{and}\quad
\widehat{\nabla_{\btheta}\rho_{\varepsilon}}(\btheta)\to \nabla_{\btheta}\rho_{P,\varepsilon}(\btheta)
\]
in probability.
Consequently,
\begin{equation}\label{eq:consistency-ghat}
\hat{\bg}_{\varepsilon}(\btheta)\;\to\;\nabla_{\btheta}J_{P,\varepsilon}(\btheta)
\qquad \text{in probability}.
\end{equation}
\end{theorem}

\begin{proof}
\textbf{Step 1: uniform convergence of $\hat F_{\btheta,\varepsilon}$.}
For fixed $\btheta$ and $\varepsilon$, the class of functions
$r\mapsto \Phi((r-S_i(\btheta))/\varepsilon)$ is uniformly bounded in $[0,1]$ and is Lipschitz in $r$ with constant $\|\varphi\|_{\infty}/\varepsilon$.
By the law of large numbers and standard stochastic equicontinuity arguments for one-dimensional Lipschitz classes (or by a DKW-type bound for smoothed indicators),
\[
\sup_{r\in \R}\abs{\hat F_{\btheta,\varepsilon}(r)-F_{\btheta,\varepsilon}(r)}\to 0
\quad\text{in probability.}
\]

\textbf{Step 2: consistency of $\hat\rho_{\varepsilon}(\btheta)$.}
Since $F_{\btheta,\varepsilon}(r)$ is nondecreasing and continuous in $r$, and by Assumption~\ref{asm:smoothed-quantile-reg} it crosses $\tau$ at a point
$\rho_{P,\varepsilon}(\btheta)$ with strictly positive derivative $\partial_r F_{\btheta,\varepsilon}$.
The uniform convergence in Step 1 implies that the level set $\{r:\hat F_{\btheta,\varepsilon}(r)\ge \tau\}$ converges to the population level set,
hence the left-quantile $\hat\rho_{\varepsilon}(\btheta)$ converges in probability to $\rho_{P,\varepsilon}(\btheta)$.

\textbf{Step 3: consistency of $\widehat{\nabla_{\btheta}\rho_{\varepsilon}}(\btheta)$.}
Write the population derivative from Lemma~\ref{lem:smoothed-quantile-deriv} as
\[
\nabla_{\btheta}\rho_{P,\varepsilon}(\btheta)
=
\frac{\E[\varphi((\rho_{P,\varepsilon}-S_{\btheta})/\varepsilon)\,\nabla_{\btheta}s_{\btheta}(U)]}
{\E[\varphi((\rho_{P,\varepsilon}-S_{\btheta})/\varepsilon)]}.
\]
The numerator and denominator are expectations of integrable random variables by Assumptions~\ref{asm:smoothing} and~\ref{asm:score-diff},
so their empirical averages converge in probability to the expectations.
Since $\hat\rho_{\varepsilon}(\btheta)\to \rho_{P,\varepsilon}(\btheta)$ and $\varphi$ is continuous and bounded,
a continuous mapping argument yields
\[
\sum_{i=1}^{n_{\mathrm{tune}}}\omega_i(\btheta)\Big/ n_{\mathrm{tune}}
\to
\E\left[\varphi\left(\frac{\rho_{P,\varepsilon}(\btheta)-S_{\btheta}}{\varepsilon}\right)\right],
\]
and similarly for the weighted average of $\nabla_{\btheta}s_{\btheta}(U_i)$.
Because the limit denominator is strictly positive by Assumption~\ref{asm:smoothed-quantile-reg} and~\eqref{eq:dFr},
the ratio converges to $\nabla_{\btheta}\rho_{P,\varepsilon}(\btheta)$.

\textbf{Step 4: consistency of $\hat{\bg}_{\varepsilon}(\btheta)$.}
By Step 2, $(\btheta,\hat\rho_{\varepsilon}(\btheta))\to (\btheta,\rho_{P,\varepsilon}(\btheta))$ in probability.
Assumptions~\ref{asm:dual-unique} and~\ref{asm:sensitivity-cont} imply that all envelope terms in~\eqref{eq:ghat} vary continuously in $(\btheta,\rho)$.
Thus the first two factors in~\eqref{eq:ghat} converge to their population counterparts at $(\btheta,\rho_{P,\varepsilon}(\btheta))$.
Combining with Step 3 gives convergence of $\hat{\bg}_{\varepsilon}(\btheta)$ to the chain-rule gradient $\nabla_{\btheta}J_{P,\varepsilon}(\btheta)$.
\end{proof}

\subsection{Smoothing bias: connecting $J_{P,\varepsilon}$ to $J_P$ as $\varepsilon\to 0$}

Theorem~\ref{thm:consistency-eps} shows that for fixed $\varepsilon>0$, the tuning-based gradient estimator targets the gradient of the smoothed profiled objective $J_{P,\varepsilon}$.
We now state conditions under which $J_{P,\varepsilon}$ (and its gradient) approach the oracle objective $J_P$ as $\varepsilon\downarrow 0$.

\begin{assumption}[Local smoothness of the score CDF]\label{asm:cdf-smooth}
Fix $\btheta$ and let $\rho_P(\btheta)$ be the $\tau$-quantile.
Assume $F_{\btheta}$ is continuously differentiable in a neighborhood of $\rho_P(\btheta)$ with density $f_{\btheta}$ satisfying
$f_{\btheta}(\rho_P(\btheta))\ge f_{\min}>0$, and $f_{\btheta}$ is Lipschitz on that neighborhood.
\end{assumption}

\begin{lemma}[Smoothed quantile converges to the true quantile]\label{lem:rhoeps-to-rho}
Assume Assumptions~\ref{asm:smoothing} and~\ref{asm:cdf-smooth}. Fix $\btheta$.
Then $\rho_{P,\varepsilon}(\btheta)\to \rho_P(\btheta)$ as $\varepsilon\downarrow 0$.
Moreover, there exists $C<\infty$ (depending on $f_{\min}$ and local smoothness) such that for sufficiently small $\varepsilon$,
\[
\abs{\rho_{P,\varepsilon}(\btheta)-\rho_P(\btheta)}\le C\,\varepsilon.
\]
\end{lemma}

\begin{proof}
The smoothed CDF $F_{\btheta,\varepsilon}$ is a mollification of the (possibly nonsmooth) indicator CDF.
Under Assumption~\ref{asm:smoothing}, $\Phi((r-S)/\varepsilon)$ converges pointwise to $\mathbf{1}\{S\le r\}$ as $\varepsilon\downarrow 0$,
and is dominated by $1$, so $F_{\btheta,\varepsilon}(r)\to F_{\btheta}(r)$ for each $r$ by dominated convergence.

Under Assumption~\ref{asm:cdf-smooth}, $F_{\btheta}$ is strictly increasing near $\rho_P(\btheta)$ with slope bounded below by $f_{\min}$.
Since $F_{\btheta,\varepsilon}$ converges uniformly to $F_{\btheta}$ on compact neighborhoods of $\rho_P(\btheta)$ (by standard properties of mollifiers under local Lipschitzness),
the solutions to $F_{\btheta,\varepsilon}(r)=\tau$ converge to the solution to $F_{\btheta}(r)=\tau$, i.e., $\rho_{P,\varepsilon}(\btheta)\to \rho_P(\btheta)$.
The linear bound $\abs{\rho_{P,\varepsilon}-\rho_P}\le C\varepsilon$ follows from a first-order Taylor expansion of $F_{\btheta}$ around $\rho_P(\btheta)$,
using the lower bound $f_{\min}$ to invert the CDF and the fact that the mollification error in CDF space is $O(\varepsilon)$ under local smoothness.
\end{proof}

\begin{remark}[Gradient convergence as $\varepsilon\downarrow 0$]
Under additional regularity ensuring that the implicit derivative formula~\eqref{eq:quantile-ift} is valid and stable (Assumption~\ref{asm:quantile-reg}),
and that $\btheta\mapsto V(\btheta,\rho)$ is smooth enough in a neighborhood of $\rho_P(\btheta)$,
one can show $\nabla_{\btheta}J_{P,\varepsilon}(\btheta)\to \nabla_{\btheta}J_P(\btheta)$ as $\varepsilon\downarrow 0$.
The key idea is: $\rho_{P,\varepsilon}(\btheta)\to \rho_P(\btheta)$ (Lemma~\ref{lem:rhoeps-to-rho}),
$\partial_\rho V$ is continuous in $\rho$ (by Assumption~\ref{asm:sensitivity-cont} and dual continuity),
and $\nabla_{\btheta}\rho_{P,\varepsilon}(\btheta)$ approximates the boundary conditional expectation defining $\nabla_{\btheta}\rho_P(\btheta)$.
We omit the most technical measure-theoretic steps here; in applications, the smoothing is used explicitly, so $J_{P,\varepsilon}$ is the effective target.
\end{remark}

\subsection{Putting it together: tuned gradients + final conformal calibration}

The final deployed radius is \emph{not} $\hat\rho_{\varepsilon}$ (tuning), but the conformal radius $\hat\rho_{\tau}$ computed on $\cD_{\mathrm{cal}}$.
This preserves finite-sample marginal coverage regardless of tuning approximations.

\begin{theorem}[Coverage is preserved under tuned training]\label{thm:coverage-preserved}
Assume Assumption~\ref{asm:exchange} holds for the calibration sample $\cD_{\mathrm{cal}}$ and a fresh $U_{\mathrm{new}}$.
Let $\hat\btheta$ be any (possibly random) shape computed using $\cD_{\mathrm{tr}}$ and $\cD_{\mathrm{tune}}$ only (no dependence on $\cD_{\mathrm{cal}}$).
Compute $\hat\rho_\tau$ on $\cD_{\mathrm{cal}}$ using~\eqref{eq:conformal-rho} with shape fixed to $\hat\btheta$.
Then
\[
\Prob\big(U_{\mathrm{new}}\in \cU_{\hat\btheta,\hat\rho_\tau}\big)\ge \tau.
\]
\end{theorem}

\begin{proof}
Condition on $\hat\btheta$. By construction, $\hat\btheta$ is independent of $\cD_{\mathrm{cal}}$.
Applying Theorem~\ref{thm:conformal} conditional on $\hat\btheta$ yields
$\Prob(U_{\mathrm{new}}\in \cU_{\hat\btheta,\hat\rho_\tau}\mid \hat\btheta)\ge \tau$ almost surely.
Taking expectations over $\hat\btheta$ gives the unconditional bound.
\end{proof}

\section{Objective degradation from calibration (optional quantitative bounds)}

This section provides a standard way to translate radius mismatch into robust objective degradation.
It is optional for the core logical justification of Option 2 but useful in a full theory write-up.

\begin{assumption}[Lipschitzness in radius]\label{asm:lipschitz-rho}
There exists $L_{\rho}<\infty$ such that for all $\btheta$ and all $\rho_1,\rho_2>0$,
\[
\abs{V(\btheta,\rho_2)-V(\btheta,\rho_1)} \le L_{\rho}\abs{\rho_2-\rho_1}.
\]
\end{assumption}

\begin{proposition}[Radius error implies objective error]\label{prop:obj-gap}
Assume Assumption~\ref{asm:lipschitz-rho}.
Fix $\btheta$. Then for any (possibly random) radius $\hat\rho$,
\[
\abs{V(\btheta,\hat\rho)-V(\btheta,\rho_P(\btheta))}
\le
L_{\rho}\,\abs{\hat\rho-\rho_P(\btheta)}.
\]
\end{proposition}

\begin{proof}
Apply Assumption~\ref{asm:lipschitz-rho} with $\rho_1=\rho_P(\btheta)$ and $\rho_2=\hat\rho$.
\end{proof}

For i.i.d.\ data and continuous score distributions, conformal order statistics admit an exact Beta distribution after applying the probability integral transform.

\begin{lemma}[Order statistic in CDF space]\label{lem:beta}
Assume $U_1,\dots,U_{n_{\mathrm{cal}}}$ are i.i.d.\ and $S_i=s_{\btheta}(U_i)$ has continuous CDF $F_{\btheta}$.
Let $S_{(k)}$ be the $k$-th order statistic. Then
\[
F_{\btheta}\big(S_{(k)}\big)\;\sim\;\mathrm{Beta}\big(k,\;n_{\mathrm{cal}}+1-k\big).
\]
\end{lemma}

\begin{proof}
By the probability integral transform, $T_i:=F_{\btheta}(S_i)$ are i.i.d.\ uniform on $[0,1]$.
Order statistics of i.i.d.\ uniforms satisfy $T_{(k)}\sim \mathrm{Beta}(k,n_{\mathrm{cal}}+1-k)$.
Finally, $T_{(k)}=F_{\btheta}(S_{(k)})$ almost surely because $F_{\btheta}$ is continuous and strictly increasing on the support.
\end{proof}

Combining Lemma~\ref{lem:beta} with a density lower bound yields a quantitative radius error rate; then Proposition~\ref{prop:obj-gap} yields objective degradation rates.
We omit the final algebraic corollaries here, since the exact form depends on the desired tail probability and local density constants.

\section{Algorithm (Option 2)}

\begin{algorithm}[t]
\caption{Option 2: Tuned Profiled-Gradient Training + Final Conformal Calibration}
\label{alg:opt2}
\begin{algorithmic}[1]
\STATE \textbf{Input:} training split $\cD_{\mathrm{tr}}$, tuning split $\cD_{\mathrm{tune}}$, calibration split $\cD_{\mathrm{cal}}$, target $\tau$, smoothing $\varepsilon$, steps $K$, step sizes $\{\eta_k\}$
\STATE \textbf{Initialize:} shape parameter $\btheta_0\in \Theta$
\FOR{$k=0,1,\dots,K-1$}
    \STATE \textbf{(Tuning)} Compute $\hat\rho_{\varepsilon}(\btheta_k)$ from $\cD_{\mathrm{tune}}$ via~\eqref{eq:rhohat-eps}
    \STATE \textbf{(Tuning)} Compute $\widehat{\nabla_{\btheta}\rho_{\varepsilon}}(\btheta_k)$ via~\eqref{eq:rhohat-grad}
    \STATE \textbf{(Robust solve)} Solve the robust problem~\eqref{eq:robust} at $(\btheta_k,\hat\rho_{\varepsilon}(\btheta_k))$ to obtain unique $\bmu^*(\btheta_k,\hat\rho_{\varepsilon}(\btheta_k))$
    \STATE \textbf{(Gradient)} Form $\hat{\bg}_{\varepsilon}(\btheta_k)$ using~\eqref{eq:ghat}
    \STATE \textbf{(Update)} $\btheta_{k+1}\leftarrow \Pi_{\Theta}\big(\btheta_k - \eta_k\,\hat{\bg}_{\varepsilon}(\btheta_k)\big)$
\ENDFOR
\STATE \textbf{Output shape:} $\hat\btheta:=\btheta_K$
\STATE \textbf{(Final conformal calibration)} Compute $\hat\rho_\tau$ on $\cD_{\mathrm{cal}}$ using~\eqref{eq:conformal-rho} with shape $\hat\btheta$
\STATE \textbf{Deploy} uncertainty set $\cU_{\hat\btheta,\hat\rho_\tau}$
\end{algorithmic}
\end{algorithm}

\begin{remark}[Why this is ``one step'' when $P$ is known]
When $P$ is known, $\rho_P(\btheta)$ is a deterministic function and $J_P(\btheta)=V(\btheta,\rho_P(\btheta))$ can be optimized directly
using the exact gradient~\eqref{eq:true-grad-expanded}. In the distribution-free regime, the tuning split provides a statistically consistent approximation of the missing term,
and the independent calibration split guarantees exact finite-sample marginal coverage for deployment.
\end{remark}

\end{document}
